\section{Diseño del controlador y estimadores}
\subsection{Realimentación de estados estimados}
La ley de control para ambos esquemas será
\begin{equation}\label{eq:control}
 u[k]=-K\widehat{\mathbf{x}}[k]+N_r r[k].
\end{equation}
Los polos del controlador se asignarán mediante \(K=\operatorname{place}
(A,B,p_c)\), verificando los valores propios de \(A-BK\). En el caso
SISO, \(D=0\) y seguimiento de referencia constante, se utilizará
\begin{equation}
 N_r=\left[C(I-A+BK)^{-1}B\right]^{-1}.
\end{equation}
Los cálculos no suponen que el comportamiento con saturación sea exactamente
el del modelo lineal sin restricciones.

\subsection{Estimador de predicción}
La innovación se calcula con la salida disponible \(y[k]\), pero la
estimación corregida se utiliza en el siguiente instante:
\begin{equation}\label{eq:predictor}
\widehat{\mathbf{x}}[k+1]=A\widehat{\mathbf{x}}[k]+Bu[k]
 +L_p\big(y[k]-C\widehat{\mathbf{x}}[k]\big).
\end{equation}
Con \(D=0\), la dinámica de error es
\begin{equation}
 \mathbf e[k+1]=(A-L_pC)\mathbf e[k].
\end{equation}
Por dualidad, \(L_p=\operatorname{place}(A^T,C^T,p_e)^T\).

\subsection{Estimador de estado actual}
El observador actual distingue predicción y corrección en la misma muestra:
\begin{align}
\mathbf z[k]&=A\widehat{\mathbf{x}}[k-1]+Bu[k-1],\\
\widehat{\mathbf{x}}[k]&=\mathbf z[k]+L_a\big(y[k]-C\mathbf z[k]\big).
\end{align}
La ley de control~\eqref{eq:control} utiliza la estimación corregida.
La matriz de transición del error corregido es
\begin{equation}
 (I-L_a C)A=A-L_a C A.
\end{equation}
En consecuencia, para asignar sus polos se empleará
\(L_a=\operatorname{place}(A^T,(CA)^T,p_e)^T\), comprobando previamente
observabilidad y rango adecuados. Se evitará equiparar directamente
\(L_p\) con \(L_a\), pues representan arquitecturas diferentes.

\subsection{Experimentos de ubicación de polos}
Se mantendrá el mismo \(K\) en las cuatro comparaciones:
(predicción, E1), (predicción, E2), (actual, E1) y (actual, E2).
Las dos parejas del observador deberán situarse dentro del círculo unidad
y responder a un criterio de velocidad y sensibilidad al ruido.
La elección definitiva de polos y los parámetros numéricos se registrarán
en las tablas de diseño, no se fijan en esta versión inicial.
